\chapter{Reference}\label{ch:reference}

\section{The \texttt{ATE} table}

\begin{center}
\footnotesize
\begin{tabular}{@{}>{\raggedright\arraybackslash}p{5.6cm}>{\raggedright\arraybackslash}p{9.4cm}@{}}
\toprule
Call & Returns / does \\
\midrule
\texttt{ATE.open(deviceType, params [, onError])} & opens a Device, returns an opaque handle \\
\texttt{ATE.read(dev [, params] [, onError])} & \texttt{\{data, count, partial\}}, or
  \texttt{err} on failure \\
\texttt{ATE.write(dev, params [, onError])} & \texttt{\{ok, count\}}; \texttt{count} includes
  the terminator \\
\texttt{ATE.close(dev [, params] [, onError])} & \texttt{\{ok = true\}} \\
\texttt{ATE.flush(dev)} & \texttt{\{ok = true, count = n\}} bytes discarded \\
\texttt{ATE.sleep(seconds)} & pauses the step; an abort wakes it at once \\
\texttt{ATE.input(n)} & the current step's $n$-th Step Input \\
\texttt{ATE.report(passFail, msg, actual, expected)} & records the step's result; returns
  \texttt{true} \\
\texttt{ATE.log(level, \dots)} & writes a line to the Test Run's log \\
\texttt{ATE.addEquipment(name, serial, connection, calDue, notes)} & declares Equipment for the
  Test Report \\
\texttt{ATE.run} & read-only table describing the current Test Run \\
\bottomrule
\end{tabular}
\end{center}

\section{The enums}

The enum tables are globals, so a schedule reads them directly:

\begingroup\footnotesize
\begin{verbatim}
Result:      PASS, FAIL, ERROR, SKIPPED, NOTRUN, FORCED_NOTRUN, NOREPORT
Log:         DEBUG, INFO, WARNING, ERROR
DeviceType:  TCP, SERIAL, UDP, VISA
Parity:      NONE, EVEN, ODD, MARK, SPACE
FlowControl: NONE, RTSCTS, XONXOFF, DSRDTR
\end{verbatim}
\endgroup

\section{\texttt{ScheduleParams} and \texttt{StepParams}}

The setters are what a schedule calls; each one returns the object, so they chain:

\begingroup\footnotesize
\begin{verbatim}
ScheduleParams(name, description): setDocTitle, setDocId, setDocRev, setPreTestAction,
  setPostTestAction, setStopOnFail, setSchedule
StepParams(name): setStepNo, setEnable, setInputs, setAction, setInclude, setPrefix
\end{verbatim}
\endgroup

The builders in \texttt{ate.devices} have their own setters: \texttt{setHost},
\texttt{setPort}, \texttt{setLocalPort}, \texttt{setResource}, \texttt{setBaud},
\texttt{setDataBits}, \texttt{setStopBits}, \texttt{setParity}, \texttt{setFlowControl},
\texttt{setTimeout}, \texttt{setReadTerm}, \texttt{setWriteTerm}, \texttt{setSize} and
\texttt{Write(data)}.

\section{What the Test Report shows, and where it goes}

ATE writes the report when the run ends, without asking:

\begingroup\footnotesize
\begin{verbatim}
<UUT Type>/<UUT Type>_<serial>_<yyyymmdd-HHMMSS>_<PASS|FAIL>.html
\end{verbatim}
\endgroup

under \texttt{reports.directory}, and the same name with \texttt{.pdf} when the PDF is
rendered: an elevated Station is asked before it writes one, and the answer defaults to Skip.
The sections are Cover,
Steps, Equipment, Configuration and Run log. The Run log in the report carries INFO, WARNING
and ERROR lines only; DEBUG lines and Device trace stay in the database.

\section{Reading a run back}

From the History screen, or from the database. Every Test Run has a UUIDv7 assigned at run
start, and its rows are the run itself, its Media Releases, its files, its schedule instances,
its nodes, its Actions, its Attempts, its Operator inputs, its Equipment and its log lines. The
End Reason tells you how the run finished: \texttt{COMPLETED}, \texttt{PRETEST\_FAILED},
\texttt{STOPPED\_ON\_FAIL}, \texttt{ABORTED}, or \texttt{INTERRUPTED} when ATE itself stopped
before the run could.

\cando{look up any part of a finished Test Run, and name the enum value or record field you are
reading.}
